% =====================================================================
% 3. EXPLORATORY DATA ANALYSIS (EDA)
% =====================================================================
\section{Exploratory Data Analysis (EDA)}

Notebook \texttt{02\_eda} follows the sections of the course EDA notebook (quality, target and
leakage, univariate analysis, normality, tests and effect sizes, outliers, correlation and VIF,
relevance, conditional probability, multivariate view, baseline) and adds an EDA of the regression
target. Notebook \texttt{03\_unsupervised} contains the PCA, MDS/t-SNE and clustering. The statistics
are implemented in \texttt{src/eda.py} and \texttt{src/\allowbreak{}unsupervised.\allowbreak{}py}; this section presents the
graphical results, and their interpretation is in the PRDL report.

\subsection*{Target, age and leakage}

The classes are moderately imbalanced (\NLGG{} LGG, \NGBM{} GBM; majority-class accuracy \MajorityAcc).
Age differs strongly between grades (Figure~\ref{fig:mllb-age}): the empirical cumulative distribution
functions are separated by a Kolmogorov--Smirnov distance of \AgeKS, and the best single threshold
(Youden's index) is \YoudenAge{} years, with sensitivity \YoudenSens{} and specificity \YoudenSpec{} on
the same data (optimistic by construction). Q--Q plots and Shapiro--Wilk tests (in the notebook) show
non-normal age distributions, which is why the Mann--Whitney test is used.

\begin{figure}[H]
\centering
\includegraphics[width=\textwidth]{eda_age_by_grade.png}
\caption{Age at diagnosis by grade: histogram, box and violin plot, and ECDF.}
\label{fig:mllb-age}
\end{figure}

\subsection*{Relevance and redundancy of the predictors}

Figure~\ref{fig:mllb-relevance} ranks every predictor by its AUC used alone, with 2000-resample
stratified bootstrap intervals (\texttt{evaluation.\allowbreak{}bootstrap\_\allowbreak{}ci} with a rank-based AUC written in
NumPy so that it can be resampled quickly). The correlation matrix (Figure~\ref{fig:mllb-corr}) is
drawn with a diverging blue--grey--red scale centred at zero; the maximum VIF is \MaxVIFEDA.
Figure~\ref{fig:mllb-condprob} shows the empirical $P(\mathrm{GBM})$ by age quintile and with and
without each mutation, with Wilson 95\% intervals.

\begin{figure}[H]
\centering
\begin{minipage}[t]{0.49\textwidth}\centering
\includegraphics[width=\textwidth]{eda_feature_relevance.png}
\caption{Single-feature AUC with 95\% bootstrap intervals (colour: direction of the effect).}
\label{fig:mllb-relevance}
\end{minipage}\hfill
\begin{minipage}[t]{0.49\textwidth}\centering
\includegraphics[width=\textwidth]{eda_feature_correlation.png}
\caption{Spearman correlation between predictors (phi coefficient for binary pairs).}
\label{fig:mllb-corr}
\end{minipage}
\end{figure}

\begin{figure}[H]
\centering
\includegraphics[width=\textwidth]{eda_conditional_probability.png}
\caption{Empirical probability of GBM by age quintile and with/without each mutation (95\% Wilson
intervals; dashed line: overall proportion of GBM).}
\label{fig:mllb-condprob}
\end{figure}

\subsection*{Quick baseline and leakage check}

A logistic regression in a \texttt{Pipeline} (scaling and one-hot encoding inside the folds), with the
same stratified 10-fold split used later, checks how much signal the EDA suggests
(Table~\ref{tab:mllb-baseline}). Adding the one-hot diagnosis gives a perfect but fake result.

\begin{table}[H]
\centering
\caption{Quick cross-validated baseline (stratified 10-fold CV, seed 42).}
\label{tab:mllb-baseline}
\resizebox{0.9\textwidth}{!}{\input{tables/eda_baseline.tex}}
\end{table}

\subsection*{Regression target}

Figure~\ref{fig:mllb-agetarget} shows the distribution of age and its relation with the eight
predictors with the strongest Mann--Whitney effects (\NAgeSigBH{} of \NAgeTests{} are significant after
Benjamini--Hochberg correction).

\begin{figure}[H]
\centering
\includegraphics[width=\textwidth]{eda_age_regression_target.png}
\caption{Distribution of age at diagnosis and age with/without the eight most related predictors.}
\label{fig:mllb-agetarget}
\end{figure}

\subsection*{Unsupervised analysis of the mutation profiles}

PCA was computed on centred and on standardised genes (Figure~\ref{fig:mllb-pca}); the biplot arrows
use a single scale factor for both axes so that their directions are not distorted. MDS and t-SNE use
a precomputed Jaccard distance matrix (Figure~\ref{fig:mllb-mds}). The number of clusters was chosen
with the criteria of Figure~\ref{fig:mllb-k}; the rule fixed in advance (highest $k$-means--Ward
agreement) gives $k=\KMain$. Table~\ref{tab:mllb-linkage} and Figure~\ref{fig:mllb-dendro} document why
Jaccard-based hierarchical clustering was not used for the final comparison, and
Figure~\ref{fig:mllb-clugrade} compares the clusters with the grade.

\begin{figure}[H]
\centering
\includegraphics[width=\textwidth]{unsup_pca_biplot.png}
\caption{PCA biplots of the 20 mutation indicators, centred (left) and standardised (right); points
coloured by grade after the PCA.}
\label{fig:mllb-pca}
\end{figure}

\begin{figure}[H]
\centering
\includegraphics[width=\textwidth]{unsup_mds_tsne_grade.png}
\caption{Metric MDS (stress-1 = \MDSStress) and t-SNE maps with Jaccard distance, coloured by grade.}
\label{fig:mllb-mds}
\end{figure}

\begin{figure}[H]
\centering
\includegraphics[width=\textwidth]{unsup_k_selection.png}
\caption{Choice of the number of clusters: within-cluster variation, silhouette, and agreement (ARI)
between $k$-means and Ward.}
\label{fig:mllb-k}
\end{figure}

\begin{table}[H]
\centering
\caption{Hierarchical clustering cut into $k=\KMain$ clusters for each distance and linkage.}
\label{tab:mllb-linkage}
\resizebox{\textwidth}{!}{\input{tables/unsup_linkage_comparison.tex}}
\end{table}

\begin{figure}[H]
\centering
\includegraphics[width=\textwidth]{unsup_dendrograms.png}
\caption{Dendrograms (last 30 merges) with Jaccard distance and average linkage (left) and Euclidean
distance and Ward linkage (right).}
\label{fig:mllb-dendro}
\end{figure}

\begin{figure}[H]
\centering
\includegraphics[width=\textwidth]{unsup_clusters_vs_grade.png}
\caption{Grade within each $k$-means cluster for $k=\KMain$ (main) and $k=\KSecond$ (exploratory).}
\label{fig:mllb-clugrade}
\end{figure}
